% Generated draft: requires mathematical and source review.
\documentclass[11pt]{article}
\usepackage[margin=25mm]{geometry}
\usepackage{fontspec}
\setmainfont{Latin Modern Roman}
\usepackage{amsmath,amssymb,mathtools,unicode-math}
\setmathfont{Latin Modern Math}
\usepackage{enumitem,needspace,xurl,hyperref}
\setlength{\parindent}{0pt}
\setlength{\parskip}{5pt}
\setlength{\emergencystretch}{3em}
\newcommand{\N}{\mathbb{N}}
\newcommand{\Z}{\mathbb{Z}}
\newcommand{\Q}{\mathbb{Q}}
\newcommand{\R}{\mathbb{R}}
\newcommand{\C}{\mathbb{C}}
\newcommand{\F}{\mathbb{F}}
\newcommand{\A}{\mathbb{A}}
\newcommand{\PP}{\mathbb{P}}
\newcommand{\E}{\mathbb{E}}
\newcommand{\eps}{\varepsilon}
\newcommand{\dd}{\,\mathrm{d}}
\begin{document}
\section*{UnsolvedMath ID: 10900024}
\textbf{Source problem code:} catalogue entry
\subsection{Definitions and mathematical statement}
Fix a closed surface $S$ and its group $\Gamma=\pi_1(S)$. An affine action on $\R^6$ is a homomorphism $\rho:\Gamma\to\operatorname{Aff}(\R^6)$, expressible as
\[
 \rho(\gamma)x=L(\gamma)x+u(\gamma),\qquad
 L:\Gamma\to\operatorname{GL}_6(\R),\qquad
 u(\gamma\delta)=u(\gamma)+L(\gamma)u(\delta).
\]
It is properly discontinuous if, for every compact $K\subset\R^6$, the set
\[
 \{\gamma\in\Gamma:\rho(\gamma)K\cap K\ne\varnothing\}
\]
is finite.

Kassel and Mann's question, item 4.5, asks to classify all such properly discontinuous affine actions for the given surface group. A usual classification convention is affine conjugacy: $\rho'$ and $\rho$ are identified when $\rho'(\gamma)=h\rho(\gamma)h^{-1}$ for a single affine automorphism $h$ and every $\gamma$. The excerpt does not specify this equivalence or whether precomposition by automorphisms of $\Gamma$ is also allowed. It also does not specify the genus or orientability of $S$; the commonly intended closed hyperbolic surface case must not silently replace the written scope.

No cocompactness, prescribed linear representation, or Hitchin/Fuchsian restriction is imposed in the question. The registry's progress for special semisimple linear parts, including its discussion of an $\operatorname{SL}_3(\R)$-related setting, is context rather than a restriction on the requested classification in $\R^6$.
\subsection{Short English statement}
For a fixed closed surface group, describe every affine action on $\R^6$ that is properly discontinuous, with the intended notion of equivalence made explicit.
\subsection{Sources}
\begin{itemize}
\item Local registry: \url{data/unsolvedmath/problems.json}, permanent UnsolvedMath ID $10900024$, exact source code \texttt{catalogue entry}, attributed to Kassel and Mann.
\item Delp, Hoffoss and Manning, \emph{Problems In Groups, Geometry, and Three-Manifolds} (2015), subsection 4.5, PDF page 3: \url{https://arxiv.org/pdf/1512.04620#page=3}; \url{https://arxiv.org/abs/1512.04620}. Supplied access date: 2026-07-29; source-tex extraction checked against rendered PDF text by the source record.
\item Supplied progress references: Danciger--Gu\'eritaud--Kassel, ``Proper affine actions for right-angled Coxeter groups,'' \url{https://arxiv.org/abs/1804.03132}; Labourie, ``Entropy and affine actions for surface groups,'' \url{https://arxiv.org/abs/1908.00599}; ``Deformation of Fuchsian representations and proper affine actions,'' \url{https://arxiv.org/abs/2312.16655}. The registry relates the last work to Mess, Labourie--Wentworth, Potrie--Sambarino and Smilga, and reports mixed-degree deformation obstructions. Exact scope and bibliographic matching require verification.
\end{itemize}
\end{document}
